\documentclass[a4paper]{article}
\usepackage{amsmath}
\usepackage{amssymb}
\usepackage{geometry}
\usepackage{hyperref}
\usepackage{booktabs}
\usepackage{enumerate}
\usepackage{listings}
\usepackage{xcolor}
\geometry{margin=1in}

\lstdefinestyle{stata}{
    basicstyle=\ttfamily\small,
    keywordstyle=\color{blue},
    commentstyle=\color{gray},
    frame=single,
    breaklines=true,
    showstringspaces=false,
}

\begin{document}

\title{\textit{Introduction to Econometrics}\\Assignment 4}
\author{Richard Harnisch, s5238366}
\date{\today}
\maketitle

% answer environment
\makeatletter
\newcommand{\answer}[1]{%
    \par\medskip
    \noindent
    \hspace*{-\@totalleftmargin}%
    \fbox{%
        \parbox{\dimexpr\textwidth-2\fboxsep-2\fboxrule}{#1}%
    }%
    \par\medskip
}
\makeatother

\noindent\textit{TeX code available under\\ \url{https://github.com/richardharnisch/intro-to-econometrics}}

\begin{enumerate}

    \item Let $\{z_i\}$ be an i.i.d.\ sequence with $E(z_i) = \alpha$ and $\mathrm{Var}(z_i) = \sigma^2$.

          \begin{enumerate}

              \item Derive $\mathrm{Var}\!\left(\sqrt{n}\,(\bar{z}_n - \alpha)\right)$.

                    \answer{
                        Since $\bar{z}_n = \frac{1}{n}\sum_{i=1}^n z_i$, we have
                        $\mathrm{Var}(\bar{z}_n) = \frac{1}{n^2}\sum_{i=1}^n \mathrm{Var}(z_i) = \frac{\sigma^2}{n}$,
                        where the second equality uses the i.i.d.\ assumption.
                        Therefore:
                        \[
                            \mathrm{Var}\!\left(\sqrt{n}\,(\bar{z}_n - \alpha)\right)
                            = n\,\mathrm{Var}(\bar{z}_n - \alpha)
                            = n\,\mathrm{Var}(\bar{z}_n)
                            = n \cdot \frac{\sigma^2}{n} = \sigma^2.
                        \]
                    }

              \item What is the asymptotic variance of $\bar{z}_n$? How does it compare with
                    $\mathrm{Var}\!\left(\sqrt{n}\,(\bar{z}_n - \alpha)\right)$?

                    \answer{
                        By the Central Limit Theorem (CLT), under the assumption that
                        $\sigma^2 < \infty$:
                        \[
                            \sqrt{n}\,(\bar{z}_n - \alpha) \xrightarrow{d} N(0, \sigma^2).
                        \]
                        The \emph{asymptotic variance} of $\bar{z}_n$, denoted $\mathrm{Avar}(\bar{z}_n)$,
                        is defined as the variance of this limiting distribution, i.e.\
                        $\mathrm{Avar}(\bar{z}_n) = \sigma^2$.

                        Comparing with part (a): the asymptotic variance $\sigma^2$ equals
                        $\mathrm{Var}\!\left(\sqrt{n}\,(\bar{z}_n - \alpha)\right) = \sigma^2$
                        exactly, and this holds for every $n$.
                        This is not a coincidence: for the sample mean of i.i.d.\ observations,
                        the exact finite-sample variance of $\sqrt{n}\,(\bar{z}_n - \alpha)$ is
                        already constant in $n$, so no approximation is needed.
                    }

              \item Based on asymptotic theory, how would you compute the standard error
                    used to construct 95\% confidence intervals or test statistics? Why is
                    it consistent?

                    \answer{
                        The standard error of $\bar{z}_n$ is
                        $\mathrm{SE}(\bar{z}_n) = \sqrt{\sigma^2/n} = \sigma/\sqrt{n}$.
                        Since $\sigma^2$ is unknown, we replace it with the sample variance
                        \[
                            s^2 = \frac{1}{n-1}\sum_{i=1}^n (z_i - \bar{z}_n)^2,
                        \]
                        giving the estimated standard error $\widehat{\mathrm{SE}} = s/\sqrt{n}$.

                        \textbf{Consistency.} By the Law of Large Numbers (LLN) applied to the
                        i.i.d.\ sequence $(z_i - \alpha)^2$ (which requires
                        $E[(z_i-\alpha)^2] = \sigma^2 < \infty$), we have
                        $s^2 \xrightarrow{p} \sigma^2$.
                        Since $g(x) = \sqrt{x/n}$ is continuous, the Continuous Mapping Theorem
                        (CMT) gives
                        $\widehat{\mathrm{SE}} = s/\sqrt{n} \xrightarrow{p} \sigma/\sqrt{n} = \mathrm{SE}(\bar{z}_n)$,
                        so the estimated standard error is consistent.
                    }

          \end{enumerate}

    \item Let $\hat{\alpha}$ and $\hat{\beta}$ be $\sqrt{n}$-asymptotically normal estimators for
          $\alpha > 0$ and $\beta > 0$ respectively. What would be a consistent estimator for
          $\delta = \log(\alpha/\beta)$?

          \answer{
              A consistent estimator is
              \[
                  \hat{\delta} = \log\!\left(\hat{\alpha}/\hat{\beta}\right)
                  = \log\hat{\alpha} - \log\hat{\beta}.
              \]

              \textbf{Why it is consistent.}
              Since $\hat{\alpha}$ and $\hat{\beta}$ are $\sqrt{n}$-asymptotically normal, they are
              consistent: $\hat{\alpha} \xrightarrow{p} \alpha > 0$ and
              $\hat{\beta} \xrightarrow{p} \beta > 0$.
              The function $g(a, b) = \log(a/b)$ is continuous on the domain
              $\{(a,b): a > 0,\, b > 0\}$, which contains the limit point $(\alpha, \beta)$.
              By the CMT:
              \[
                  \hat{\delta} = g(\hat{\alpha}, \hat{\beta}) \xrightarrow{p}
                  g(\alpha, \beta) = \log(\alpha/\beta) = \delta.
              \]
          }

    \item Consider the model $y_i = x_i\beta + \epsilon_i$ with $E(\epsilon_i \mid x_i) = 0$,
          $x_i \in \mathbb{R}$, and i.i.d.\ observations, and the two estimators from Exercise 2
          of the Week 4 tutorial:
          \[
              \hat{\beta} = \frac{\sum_{i=1}^n x_i y_i}{\sum_{i=1}^n x_i^2}, \qquad
              \tilde{\beta} = \frac{1}{n}\sum_{i=1}^n \frac{y_i}{x_i}.
          \]

          \begin{enumerate}[(i)]
              \item \textbf{Estimator $\hat{\beta}$.}

                    \answer{
                        \textbf{Reformulation.} Substituting $y_i = x_i\beta + \epsilon_i$:
                        \[
                            \hat{\beta} = \beta
                            + \frac{\frac{1}{n}\sum_{i=1}^n x_i \epsilon_i}
                            {\frac{1}{n}\sum_{i=1}^n x_i^2},
                        \]
                        so $\sqrt{n}(\hat{\beta} - \beta)
                            = \dfrac{\frac{1}{\sqrt{n}}\sum_{i=1}^n x_i\epsilon_i}
                            {\frac{1}{n}\sum_{i=1}^n x_i^2}$.

                        \textbf{Additional moment condition.}
                        The consistency proof used $E[x_i^2] < \infty$ (LLN on denominator) and
                        $E[|x_i \epsilon_i|] < \infty$ (LLN on numerator).
                        To apply the CLT to $\{x_i \epsilon_i\}$, we additionally require
                        \[
                            E\!\left[x_i^2\,\epsilon_i^2\right] < \infty.
                        \]
                        This ensures $\mathrm{Var}(x_i\epsilon_i) = E[x_i^2\epsilon_i^2] - (E[x_i\epsilon_i])^2
                            = E[x_i^2\epsilon_i^2]$ is finite (using $E[x_i\epsilon_i]=0$).

                        \textbf{Asymptotic distribution.}
                        Denote $\mu_2 = E[x_i^2] \neq 0$ and
                        $\sigma_{x\epsilon}^2 = E[x_i^2\epsilon_i^2]$.
                        By the CLT:
                        \[
                            \frac{1}{\sqrt{n}}\sum_{i=1}^n x_i\epsilon_i
                            \xrightarrow{d} N\!\left(0,\, \sigma_{x\epsilon}^2\right).
                        \]
                        Since $\frac{1}{n}\sum x_i^2 \xrightarrow{p} \mu_2$ by the LLN,
                        Slutsky's theorem gives:
                        \[
                            \sqrt{n}\,(\hat{\beta} - \beta)
                            \xrightarrow{d} N\!\!\left(0,\;
                            \frac{E\!\left[x_i^2\epsilon_i^2\right]}{\left(E\!\left[x_i^2\right]\right)^2}
                            \right).
                        \]
                    }

              \item \textbf{Estimator $\tilde{\beta}$.}

                    \answer{
                    \textbf{Reformulation.} Substituting $y_i = x_i\beta + \epsilon_i$:
                    \[
                        \tilde{\beta} = \beta + \frac{1}{n}\sum_{i=1}^n \frac{\epsilon_i}{x_i},
                    \]
                    so $\sqrt{n}(\tilde{\beta} - \beta) = \dfrac{1}{\sqrt{n}}\sum_{i=1}^n
                        \dfrac{\epsilon_i}{x_i}$.

                    \textbf{Additional moment condition.}
                    The consistency proof used $P(x_i = 0) = 0$ and
                    $E[|\epsilon_i/x_i|] < \infty$ (LLN on the sum). To apply the CLT to
                    $\{\epsilon_i/x_i\}$, we additionally require
                    \[
                        E\!\left[\frac{\epsilon_i^2}{x_i^2}\right] < \infty.
                    \]
                    This ensures $\mathrm{Var}(\epsilon_i/x_i) = E[\epsilon_i^2/x_i^2]$
                    is finite (using $E[\epsilon_i/x_i] = E[(1/x_i)\,E[\epsilon_i \mid x_i]] = 0$
                    by the Law of Iterated Expectations).

                    \textbf{Asymptotic distribution.}
                    By the CLT:
                    \[
                        \sqrt{n}\,(\tilde{\beta} - \beta)
                        \xrightarrow{d} N\!\!\left(0,\;
                        E\!\left[\frac{\epsilon_i^2}{x_i^2}\right]
                        \right).
                    \]
                    }

          \end{enumerate}

    \item Consider the OLS regression:
          \[
              \log \mathrm{wage}_i = \beta_1\,\mathrm{education}_i
              + \beta_2\,\mathrm{experience}_i
              + \beta_3\,\frac{\mathrm{experience}_i^2}{100}
              + \beta_4 + \varepsilon_i.
          \]

          \begin{enumerate}

              \item How would you interpret $\beta_1$?

                    \answer{
                        Since the dependent variable is $\log\mathrm{wage}$, $\beta_1$ is the
                        causal effect of one additional year of education on log wages, holding
                        experience constant.
                        In percentage terms, an additional year of education is associated with
                        an approximate $100\beta_1\%$ increase in wages.
                        With the estimated value $b_1 = 0.118$, one extra year of education
                        raises wages by approximately $11.8\%$.
                    }

              \item What is the marginal impact of experience on log wages?

                    \answer{
                        Differentiating the model with respect to experience:
                        \[
                            \frac{\partial \log\mathrm{wage}_i}{\partial\,\mathrm{experience}_i}
                            = \beta_2 + \frac{2\beta_3}{100}\,\mathrm{experience}_i
                            = \beta_2 + \frac{\beta_3}{50}\,\mathrm{experience}_i.
                        \]
                        This is a non-constant marginal effect that depends on the level of
                        experience. With $b_2 = 0.016 > 0$ and $b_3 = -0.022 < 0$, the
                        marginal return to experience is positive but decreasing (concave in
                        experience). Multiplied by 100, it gives the approximate percentage
                        increase in wages per additional year of experience at a given level.
                    }

              \item $\theta = 100\beta_2 + 20\beta_3$ measures the percentage return to
                    experience for individuals with $X$ years. What is $X$?

                    \answer{
                        The percentage return to experience at $X$ years is
                        $100 \times \left(\beta_2 + \dfrac{2\beta_3}{100}\,X\right)
                            = 100\beta_2 + 2\beta_3\,X$.

                        Setting this equal to $\theta = 100\beta_2 + 20\beta_3$:
                        \[
                            2\beta_3\,X = 20\beta_3
                            \implies X = 10.
                        \]
                        So $\theta$ measures the percentage return to experience for individuals
                        with \textbf{10 years} of experience.
                    }

              \item Standard errors and test of returns to education.

                    \answer{
                        The standard errors are $\mathrm{SE}(b_j) = \sqrt{\hat{V}_\beta[j,j]}$
                        where $\hat{V}_\beta[j,j]$ are the diagonal entries of $\hat{V}_\beta$
                        (in units of $10^{-4}$):
                        \begin{center}
                            \begin{tabular}{lcc}
                                \toprule
                                Parameter & $\hat{V}_\beta[j,j]$   & $\mathrm{SE}(b_j)$                            \\
                                \midrule
                                $b_1$     & $0.632 \times 10^{-4}$ & $\sqrt{0.632 \times 10^{-4}} \approx 0.00795$ \\
                                $b_2$     & $0.390 \times 10^{-4}$ & $\sqrt{0.390 \times 10^{-4}} \approx 0.00625$ \\
                                $b_3$     & $1.48  \times 10^{-4}$ & $\sqrt{1.48  \times 10^{-4}} \approx 0.01217$ \\
                                $b_4$     & $246   \times 10^{-4}$ & $\sqrt{246   \times 10^{-4}} \approx 0.15684$ \\
                                \bottomrule
                            \end{tabular}
                        \end{center}

                        \textbf{Test: $H_0$: $\beta_1 = 0.10$ vs.\ $H_1$: $\beta_1 \neq 0.10$
                            at $\alpha = 5\%$.}

                        Under $H_0$, the asymptotic $t$-statistic is:
                        \[
                            t = \frac{b_1 - 0.10}{\mathrm{SE}(b_1)}
                            = \frac{0.118 - 0.10}{0.00795}
                            \approx 2.264.
                        \]
                        For large $n$, $t \xrightarrow{d} N(0,1)$ under $H_0$, so we compare
                        to the critical value $z_{0.025} \approx 1.960$.
                        Since $|t| = 2.264 > 1.960$, we \textbf{reject} $H_0$ at the 5\% level.
                        The data suggest that the returns to education are significantly
                        different from 10\%.
                    }

              \item Test whether $\theta = 100\beta_2 + 20\beta_3$ is different from 0.

                    We test $H_0: R\beta = r$ with $R = (0,\, 100,\, 20,\, 0)$ and $r = 0$
                    ($\#r = 1$ restriction). Since
                    $\hat{V}_\beta = \frac{1}{n}\widehat{\mathrm{Avar}}(b)$,
                    the test statistic simplifies to
                    \[
                        W = (Rb - r)^\prime
                        \bigl(R\hat{V}_\beta R^\prime\bigr)^{-1}
                        (Rb - r)
                        \xrightarrow{d} \chi^2(1).
                    \]

                    \textbf{Stata code:}

                    \begin{lstlisting}[style=stata]
* Coefficient vector (column)
matrix b = (0.118 \ 0.016 \ -0.022 \ 0.947)

* V_hat_beta (entries * 1e-4)
matrix Vbeta = (0.0000632,  0.0000131, -0.0000143, -0.00111   ///
               \ 0.0000131,  0.0000390, -0.0000731, -0.000625  ///
               \ -0.0000143, -0.0000731,  0.000148,  0.000943  ///
               \ -0.00111,   -0.000625,   0.000943,  0.0246)

* R row vector: theta = 100*beta2 + 20*beta3
matrix R = (0, 100, 20, 0)

* Rb (r = 0)
matrix diff = R * b
matlist diff

* R * Vbeta * R'
matrix RVR = R * Vbeta * R'
matlist RVR

* Wald statistic
matrix W = diff' * inv(RVR) * diff
display W[1,1]

* p-value and 5% critical value
display 1 - chi2(1, W[1,1])
display invchi2(1, 0.95)
                    \end{lstlisting}

                    \answer{
                        $Rb - r = 100(0.016) + 20(-0.022) = 1.16$ and
                        $R\hat{V}_\beta R' = 100^2(0.390) + 2(100)(20)(-0.731) + 20^2(1.48)
                            = 1568$ (all $\times 10^{-4}$) $= 0.1568$.
                        Thus $W = 1.16^2/0.1568 \approx 8.582$, which exceeds the 95th
                        percentile of $\chi^2(1) \approx 3.841$. We \textbf{reject} $H_0$
                        at the 5\% level: the percentage return to experience at 10 years
                        is significantly different from zero.
                    }

          \end{enumerate}

\end{enumerate}

\end{document}
